% ===============================================================
%  PART VI -- Results and Discussion (Chapters 19--20)
%  Every number in this Part comes from the simulator of Part V.
%  Plotted data: data/*.dat, written by simulator/make_report_data.py
% ===============================================================

\newtcolorbox{verdictbox}[2][memGreen]{
  colback=#1!6, colframe=#1, coltitle=white,
  colbacktitle=#1, fonttitle=\sffamily\bfseries\small,
  title={#2}, boxrule=0.6pt, arc=2pt, fontupper=\small,
  left=6pt, right=6pt, top=4pt, bottom=4pt
}

\part{Results and Discussion}
\label{part:results}

\begin{center}
\adjustbox{max width=\linewidth}{%
\begin{forest}
  part map,
  [Part VI
    [Ch.~19: Traffic {\&} Time
      [H1: $F$ and cores] [Traffic saved] [H2: bytes to time] [AMAT] [Memory bandwidth]]
    [Ch.~20: Testing the Predictor
      [Input knob vs DSR] [H3: DSR predicts the saving] [Survival fraction] [H4: when MOESI loses] [Verdicts]]
  ]
\end{forest}}
\end{center}

\begin{keybox}[What Part VI does]
Part~VI reports what the simulator measured and judges each hypothesis of
\cref{tab:hypotheses} against it. Chapter~19 covers cost: the size of the
forced-writeback fraction from 2 to 16 cores (H1), the traffic MOESI removes, and
whether removed traffic becomes saved time (H2). Chapter~20 tests the Dirty Sharing
Ratio as a predictor (H3) and looks for conditions under which MOESI does worse
than MESI (H4). Each section ends with a verdict: held, partly held or failed.
Unless stated otherwise a number is the mean over three seeds of paired runs
(\cref{eq:paired}), ``points'' means percentage points, and the experiments are
those of \cref{tab:experiments}.
\end{keybox}

% ===============================================================
%  Chapter 19 -- Traffic, Execution Time and AMAT
% ===============================================================
\chapter{Traffic, Execution Time and AMAT from 2 to 16 Cores}
\label{ch:res-traffic}
\index{AMAT}%

\begin{figure}[htbp]
\centering
\adjustbox{max width=\linewidth}{%
\begin{forest}
  chapter map,
  [Chapter 19
    [H1: $F$ and cores [Falls at every mix] [More readers per epoch]]
    [Traffic saved [Never more bytes] [Memory writes cut]]
    [H2: bytes to time [Gain follows utilisation] [Saturation limit]]
    [AMAT {\&} memory [AMAT is not time] [Memory-bound gains most]]
  ]
\end{forest}}
\caption[]{Structure of Chapter~19: the cost of forced
writebacks (H1), the traffic MOESI removes, and how much of it becomes time (H2).}
\label{fig:ch19-map}
\end{figure}

% ---------------------------------------------------------------
\section{H1: The Forced-Writeback Fraction from 2 to 16 Cores}
\label{sec:res-h1}
\sectionmap{[H1 [{F at 2, 4, 8, 16 cores}] [Falls at every sharing mix]
  [{Why: more readers per epoch}] [Verdict: held]]}
\index{MESI}\index{forced writeback}\index{dirty sharing}\index{dirty epoch}\index{sharing patterns!producer--consumer}\index{sharing patterns!migratory}%

H1 predicted that, for a fixed sharing structure, $F(W,N)$ falls as cores are added,
because each forced writeback is shared among more read misses.
\Cref{fig:res-h1}(a) plots $F$ under MESI against the core count for four mixes, from
almost no dirty sharing ($d=0.005$) to nothing else ($d=1$). Every line falls. The
fall is steepest when sharing is light: at $d=0.005$, $F$ drops from 17.5\,\% at 2
cores to 3.7\,\% at 16, almost a factor of five, while at $d=1$ it drops only from
44.9 to 39.0\,\%. \Cref{tab:res-h1} gives the numbers for workloads with
$d\ge0.1$, where dirty sharing is a large part of the traffic.

\begin{figure}[htbp]
\centering
\begin{bigfig}[1.0]
\begin{tikzpicture}
  \begin{axis}[name=a, width=7.4cm, height=5.6cm,
      title={\sffamily\small (a) Forced-writeback fraction $F$},
      xlabel={Cores $N$}, ylabel={$F$ under MESI (\%)},
      xmode=log, log basis x=2, xtick={2,4,8,16}, xticklabels={2,4,8,16},
      ymin=0, ymax=50, grid=major, grid style={coreGrey!15},
      label style={font=\small}, tick label style={font=\scriptsize},
      legend style={at={(1.08,-0.28)}, anchor=north, legend columns=4, font=\scriptsize, draw=none}]
    \addplot[very thick, tealC, mark=*] table[x=N, y=d0p005] {data/h1_F_vs_N.dat}; \addlegendentry{$d = 0.005$}
    \addplot[very thick, plumC, mark=square*] table[x=N, y=d0p02] {data/h1_F_vs_N.dat}; \addlegendentry{$d = 0.02$}
    \addplot[very thick, amberC, mark=triangle*] table[x=N, y=d0p1] {data/h1_F_vs_N.dat}; \addlegendentry{$d = 0.1$}
    \addplot[very thick, brickC, mark=diamond*] table[x=N, y=d1] {data/h1_F_vs_N.dat}; \addlegendentry{$d = 1$}
  \end{axis}
  \begin{axis}[at={(a.east)}, anchor=west, xshift=1.5cm, width=7.4cm, height=5.6cm,
      title={\sffamily\small (b) Read misses served by a modified copy},
      xlabel={Cores $N$}, ylabel={share of read misses (\%)},
      xmode=log, log basis x=2, xtick={2,4,8,16}, xticklabels={2,4,8,16},
      ymin=0, ymax=105, grid=major, grid style={coreGrey!15},
      label style={font=\small}, tick label style={font=\scriptsize}]
    \addplot[very thick, tealC, mark=*] table[x=N, y=d0p005] {data/h1_dirtyshare_vs_N.dat};
    \addplot[very thick, plumC, mark=square*] table[x=N, y=d0p02] {data/h1_dirtyshare_vs_N.dat};
    \addplot[very thick, amberC, mark=triangle*] table[x=N, y=d0p1] {data/h1_dirtyshare_vs_N.dat};
    \addplot[very thick, brickC, mark=diamond*] table[x=N, y=d1] {data/h1_dirtyshare_vs_N.dat};
  \end{axis}
\end{tikzpicture}
\end{bigfig}
\caption[H1: forced-writeback fraction against cores]{Experiment T1. (a) The
forced-writeback fraction falls with the core count at every mix. (b) With more
cores a smaller share of read misses finds the block modified, because the first
reader of each epoch has already forced the writeback that cleans it for the rest.
Means over seeds and producer write probabilities.}
\label{fig:res-h1}
\end{figure}

\begin{table}[htbp]
\centering
\caption{H1 in numbers: workloads with $d \ge 0.1$ (36 runs per core count).
``Dirty share'' is the fraction of read misses served by a modified copy.}
\label{tab:res-h1}
\small
\begin{tabular}{@{}c c c c c c c@{}}
\toprule
Cores & Mean $F$ (\%) & Min $F$ & Max $F$ & Dirty share (\%) & $\dsr$ & $\bar b$ (B) \\
\midrule
2 & 43.6 & 41.8 & 44.9 & 97.3 & 0.491 & 81.1 \\
4 & 40.0 & 36.8 & 42.9 & 84.4 & 0.454 & 81.8 \\
8 & 36.8 & 31.3 & 41.8 & 75.1 & 0.423 & 83.1 \\
16 & 33.4 & 25.4 & 41.4 & 67.5 & 0.393 & 85.3 \\
\bottomrule
\end{tabular}
\end{table}

\Cref{fig:res-h1}(b) shows the mechanism proposed in Chapter~16. With two cores, a
read miss on shared data can only find the block in the one other cache, the
writer, so 97\,\% of read misses on such data are dirty reads and nearly all of
them carry a forced writeback. With sixteen cores, the first reader in an epoch
forces the writeback and the readers after it find a clean block in memory, so only
67.5\,\% of read misses are dirty reads. The writeback stays at 72 bytes per epoch,
but it is shared among more 80-byte read misses, so its share falls, exactly as
$F=72/B(C)$ in \cref{eq:pc-model} falls with more consumers. Producers that write
more often create more epochs and so raise $F$: at 16 cores, $F$ rises from 31.6 to
35.3\,\% as the producer write probability goes from 0.15 to 0.5. The largest value
measured, 44.9\,\%, is the steady-state 72/160 of one producer--consumer or
migratory hand-off (\cref{tab:dsr-patterns}), and the measured range of 25 to
45\,\% for $d\ge0.1$ lies inside the 10 to 45\,\% derived in Chapter~10.

\begin{verdictbox}{Verdict on H1: held}
At every mix $F$ falls steadily from 2 to 16 cores, for the predicted reason. When
dirty sharing is substantial, the forced writeback is a third to nearly half of
MESI's coherence traffic (33.4 to 43.6\,\% on average).
\end{verdictbox}

% ---------------------------------------------------------------
\section{Traffic Saved by MOESI}
\label{sec:res-traffic}
\sectionmap{[Traffic [Never more bytes than MESI] [MESI never more than MSI]
  [Memory writes cut much more] [Growth with dirty sharing]]}
\index{MESI}\index{MOESI}\index{MSI}\index{Exclusive state}\index{forced writeback}\index{dirty sharing}%

\Cref{fig:res-traffic} compares the three protocols at 8 cores for four mixes, and
\cref{tab:res-traffic} averages over all core counts. In all 468 paired traffic
runs, including the 72 with caches as small as four blocks, MOESI moved no more
bytes than MESI, as \cref{prop:never-worse} requires. The saving tracks $F$
closely, within one point in every row of the table; Chapter~20 explains what this
does and does not show. MESI in turn never moved more bytes than MSI, and up to
3.5\,\% fewer (3.24\,\% on average with no dirty sharing). That is the \stE\ saving
of \cref{sec:fam-mesi}: a private block read and then written needs no upgrade in
MESI. It shrinks to zero as dirty sharing takes over, since shared blocks are never
in \stE.

\begin{figure}[htbp]
\centering
\begin{bigfig}[1.0]
\begin{tikzpicture}
  \begin{axis}[name=a, ybar, bar width=8pt, width=7.6cm, height=5.6cm,
      title={\sffamily\small (a) Interconnect bytes (MSI $=100$)},
      xlabel={dirty-sharing share $d$}, ylabel={relative bytes},
      xtick={0,1,2,3}, xticklabels={0.005, 0.02, 0.1, 0.5},
      ymin=0, ymax=118, enlarge x limits=0.18, grid=major, grid style={coreGrey!15},
      label style={font=\small}, tick label style={font=\scriptsize},
      legend style={at={(1.08,-0.28)}, anchor=north, legend columns=3, font=\scriptsize, draw=none}]
    \addplot[fill=coreGrey!45, draw=coreGrey] table[x=x, y=msi] {data/traffic_n8.dat}; \addlegendentry{MSI}
    \addplot[fill=mesiBlue!70, draw=mesiBlue] table[x=x, y=mesi] {data/traffic_n8.dat}; \addlegendentry{MESI}
    \addplot[fill=moesiOrange!75, draw=moesiOrange] table[x=x, y=moesi] {data/traffic_n8.dat}; \addlegendentry{MOESI}
  \end{axis}
  \begin{axis}[at={(a.east)}, anchor=west, xshift=1.5cm, ybar, bar width=10pt, width=7.6cm, height=5.6cm,
      title={\sffamily\small (b) Bytes written to memory (MESI $=100$)},
      xlabel={dirty-sharing share $d$}, ylabel={relative bytes},
      xtick={0,1,2,3}, xticklabels={0.005, 0.02, 0.1, 0.5},
      ymin=0, ymax=118, enlarge x limits=0.18, grid=major, grid style={coreGrey!15},
      label style={font=\small}, tick label style={font=\scriptsize},
      nodes near coords, nodes near coords style={font=\tiny, /pgf/number format/fixed, /pgf/number format/precision=0}]
    \addplot[fill=mesiBlue!70, draw=mesiBlue] table[x=x, y=memmesi] {data/traffic_n8.dat};
    \addplot[fill=moesiOrange!75, draw=moesiOrange] table[x=x, y=memmoesi] {data/traffic_n8.dat};
  \end{axis}
\end{tikzpicture}
\end{bigfig}
\caption[Traffic of the three protocols]{Experiment T1 at 8 cores. (a) MESI moves
slightly fewer bytes than MSI (the \stE\ state), and MOESI many fewer as dirty
sharing grows. (b) At $d=0.1$ MOESI writes less than a fifth of what MESI writes to
memory. MSI writes the same bytes to memory as MESI.}
\label{fig:res-traffic}
\end{figure}

\begin{table}[htbp]
\centering
\caption{Traffic results of experiment T1 for selected dirty-sharing shares,
averaged over 2 to 16 cores (36 runs per row).}
\label{tab:res-traffic}
\small
\begin{tabular}{@{}c c c c c@{}}
\toprule
 & & \multicolumn{2}{c}{MOESI saves vs MESI (\%)} & MESI saves \\
\cmidrule(lr){3-4}
$d$ & $F$ (\%) & interconnect bytes & memory writes & vs MSI (\%) \\
\midrule
0 & 0.0 & 0.0 & --- & 3.24 \\
0.001 & 1.7 & 1.6 & 3.8 & 3.07 \\
0.005 & 9.8 & 9.1 & 21.9 & 2.43 \\
0.01 & 16.0 & 15.3 & 36.7 & 1.98 \\
0.02 & 22.6 & 22.1 & 52.7 & 1.47 \\
0.05 & 30.5 & 30.2 & 71.9 & 0.86 \\
0.1 & 34.8 & 34.6 & 82.8 & 0.53 \\
0.5 & 39.7 & 39.7 & 95.3 & 0.14 \\
1 & 41.6 & 41.5 & 99.8 & 0.00 \\
\bottomrule
\end{tabular}
\end{table}

Memory sees most of the benefit. MOESI's cut in memory writes is about 2.4 times
its cut in interconnect bytes in every row. In these workloads the private data
fits in the caches and is rarely evicted, so almost every MESI memory write is a
forced writeback, and removing forced writebacks removes most memory writes:
82.8\,\% at $d=0.1$ and 99.8\,\% at $d=1$. Whether that saves time depends on
whether memory is a bottleneck (\cref{sec:res-memory}).

% ---------------------------------------------------------------
\section{H2: From Traffic to Execution Time}
\label{sec:res-h2}
\sectionmap{[H2 [Load sweep E1] [Gain follows utilisation]
  [Queueing explains it] [Saturation limit confirmed] [Verdict: held]]}
\index{MESI}\index{MOESI}\index{forced writeback}\index{bus utilisation}\index{Effect A and Effect B}%

H2 predicted that MOESI's time gain is near zero at low bus utilisation and rises
towards the ratio of bus-busy times, $C_{\mathrm{bus}}(\mathrm{MESI}) /
C_{\mathrm{bus}}(\mathrm{MOESI})$ (\cref{eq:sat-limit}), as utilisation nears
100\,\%. Experiment E1 keeps the workload fixed ($d=0.5$) and changes only how long
each core computes between memory operations, from 2000 cycles (an idle bus) to 2
(a saturated one). \Cref{fig:res-h2}(a) plots the time reduction against MESI's bus
utilisation, and \cref{tab:res-h2} gives the 8-core numbers.

\begin{figure}[htbp]
\centering
\begin{bigfig}[1.0]
\begin{tikzpicture}
  \begin{axis}[name=a, width=7.4cm, height=5.6cm,
      title={\sffamily\small (a) Time saved by MOESI},
      xlabel={MESI bus utilisation (\%)}, ylabel={execution-time reduction (\%)},
      xmin=0, xmax=100, ymin=0, ymax=31, grid=major, grid style={coreGrey!15},
      label style={font=\small}, tick label style={font=\scriptsize},
      legend style={at={(0.03,0.97)}, anchor=north west, font=\scriptsize, draw=none, fill=white, fill opacity=0.8, text opacity=1}]
    \addplot[very thick, mesiBlue, mark=*] table[x=Um, y=red] {data/e1_serial_N8.dat}; \addlegendentry{$N=8$, serial}
    \addplot[very thick, mesiBlue, dashed, mark=o] table[x=Um, y=red] {data/e1_overlap_N8.dat}; \addlegendentry{$N=8$, overlapped}
    \addplot[very thick, moesiOrange, mark=square*] table[x=Um, y=red] {data/e1_serial_N16.dat}; \addlegendentry{$N=16$, serial}
    \addplot[very thick, moesiOrange, dashed, mark=square] table[x=Um, y=red] {data/e1_overlap_N16.dat}; \addlegendentry{$N=16$, overlapped}
  \end{axis}
  \begin{axis}[at={(a.east)}, anchor=west, xshift=2.2cm, width=7.4cm, height=5.6cm,
      title={\sffamily\small (b) Waiting for the bus},
      xlabel={bus utilisation of that protocol (\%)}, ylabel={mean wait per transaction (cycles)},
      xmin=0, xmax=100, ymode=log, ymin=0.05, ymax=200, grid=major, grid style={coreGrey!15},
      label style={font=\small}, tick label style={font=\scriptsize},
      legend style={at={(0.03,0.97)}, anchor=north west, font=\scriptsize, draw=none, fill=white, fill opacity=0.8, text opacity=1}]
    \addplot[very thick, mesiBlue, mark=*] table[x=Um, y=Wm] {data/e1_serial_N8.dat}; \addlegendentry{MESI, $N=8$}
    \addplot[very thick, moesiOrange, mark=square*] table[x=Uo, y=Wo] {data/e1_serial_N8.dat}; \addlegendentry{MOESI, $N=8$}
    \addplot[thick, mesiBlue!60, dashed, mark=o] table[x=Um, y=Wm] {data/e1_serial_N16.dat}; \addlegendentry{MESI, $N=16$}
    \addplot[thick, moesiOrange!70, dashed, mark=square] table[x=Uo, y=Wo] {data/e1_serial_N16.dat}; \addlegendentry{MOESI, $N=16$}
  \end{axis}
\end{tikzpicture}
\end{bigfig}
\caption[H2: time saved against bus utilisation]{Experiment E1. (a) MOESI's time
gain is negligible on a lightly used bus and grows steeply as the bus fills.
(b) Waiting grows sharply with utilisation and both protocols lie on one curve;
MOESI sits further left because each of its dirty reads holds the bus for 12 cycles
instead of 22. Mean of three seeds, $d=0.5$.}
\label{fig:res-h2}
\end{figure}

\begin{table}[htbp]
\centering
\caption{Experiment E1 at 8 cores, serial writeback. ``Converted'' is the share of
the byte saving that appears as time saving. The last two columns compare the
speed-up with the limit of \cref{eq:sat-limit}.}
\label{tab:res-h2}
\small
\begin{tabular}{@{}r c c c c c c c@{}}
\toprule
Think & \multicolumn{2}{c}{Bus utilisation (\%)} & Bytes saved & Time saved & Converted & Speed-up & $C_{\mathrm{bus}}$ ratio \\
\cmidrule(lr){2-3}
(cycles) & MESI & MOESI & (\%) & (\%) & (\%) & & \\
\midrule
2000 & 2.5 & 1.7 & 32.6 & 0.3 & 1 & 1.003 & 1.459 \\
600 & 8.1 & 5.6 & 32.6 & 0.9 & 3 & 1.009 & 1.458 \\
300 & 15.5 & 10.8 & 32.4 & 1.8 & 6 & 1.019 & 1.454 \\
150 & 28.5 & 20.4 & 32.3 & 3.7 & 11 & 1.039 & 1.452 \\
80 & 46.2 & 34.3 & 31.9 & 6.8 & 21 & 1.073 & 1.445 \\
40 & 69.3 & 55.4 & 31.3 & 12.6 & 40 & 1.144 & 1.432 \\
20 & 87.0 & 77.1 & 30.8 & 20.7 & 67 & 1.260 & 1.422 \\
10 & 95.4 & 91.3 & 30.2 & 25.9 & 86 & 1.350 & 1.411 \\
5 & 98.0 & 96.8 & 30.0 & 28.0 & 93 & 1.388 & 1.406 \\
2 & 98.9 & 98.3 & 29.7 & 28.2 & 95 & 1.393 & 1.402 \\
\bottomrule
\end{tabular}
\end{table}

MOESI removes about 30\,\% of the bytes at every load, yet the time it saves ranges
from 0.3 to 28.2\,\%. Below 35\,\% bus utilisation it saves at most 3.7\,\% of the
time at 8 cores and 2.4\,\% at 16; above 65\,\% it saves 12.6 to 28.2\,\%. Put
another way, about 1\,\% of the saved bytes turn into saved time on an idle bus and
95\,\% on a saturated one. \Cref{fig:res-h2}(b) shows why. The mean wait for the bus
grows from a fraction of a cycle at low utilisation to tens of cycles near
100\,\%, as \cref{eq:util} suggests, and MESI and MOESI lie on the same curve, so
neither waits less at a given utilisation. MOESI wins because it does the same work
at a lower utilisation. At 8 cores with 40 cycles of computation, MESI's bus is
69\,\% busy and a transaction waits 11.0 cycles on average, while MOESI's is 55\,\%
busy and waits 4.6. On an idle bus there is no waiting to remove, and the only gain
left is the 9-cycle writeback of each dirty read (Effect~A of
\cref{fig:two-effects}), which is tiny next to the computation.

The last two columns of \cref{tab:res-h2} confirm the saturation limit of
\cref{eq:sat-limit}: a speed-up of 1.393 against a bus-busy ratio of 1.402 at 8
cores, and 1.318 against 1.319 at 16. At low load the ratio is still about 1.46,
since MOESI keeps the bus a third less busy, but the speed-up is 1.003. Whether the
MESI reader waits for the writeback (\cref{fig:sim-wbmode}) matters little: at 8
cores and moderate load MOESI saves 12.6\,\% of the time in serial mode and 11.0\,\%
in overlapped mode, and both reach 28.2\,\% at saturation. The forced writeback
costs time mainly through the bus time it takes up, not through the latency it adds
to one read. The 16-core curves rise at lighter per-core loads than the 8-core
ones, because twice the cores reach saturation at half the per-core demand
(\cref{eq:util}).

\begin{verdictbox}{Verdict on H2: held}
Saved traffic becomes saved time in proportion to how busy the bus is: about 1\,\%
of the byte saving at 2.5\,\% utilisation, 40\,\% at 69\,\% and 95\,\% at
saturation, where the speed-up reaches the bus-occupancy limit. A comparison run on
a lightly loaded system will see almost no difference in time, however many bytes
are saved.
\end{verdictbox}

% ---------------------------------------------------------------
\section{AMAT, Scaling with Cores and Memory Bandwidth}
\label{sec:res-amat}
\label{sec:res-memory}
\sectionmap{[AMAT and memory [Experiment E2] [Idle latencies match the report]
  [AMAT is not execution time] [E5: memory-bound gains most]]}
\index{MESI}\index{MOESI}\index{forced writeback}\index{Dirty Sharing Ratio}\index{dirty epoch}\index{AMAT}\index{bus utilisation}\index{owner}\index{Effect A and Effect B}%

Experiment E2 runs the same workload on 2 to 16 cores at a light and a heavy load
(\cref{fig:res-amat,tab:res-amat}). On a lightly loaded bus with two cores a dirty
read takes 61.1 cycles under MESI and 52.0 under MOESI, within a fraction of a
cycle of $T_{L1}+T_{\mathrm{coh}}=61$ and $T_{L1}+T_{\mathrm{c2c}}=52$
(\cref{tab:sim-latency}). As cores and load grow, queueing adds to both and more to
MESI: at 16 cores under heavy load a dirty read takes 125.5 cycles under MESI and
86.6 under MOESI.

\begin{figure}[htbp]
\centering
\begin{minipage}[c]{0.52\linewidth}\centering
\begin{tikzpicture}
  \begin{axis}[width=\linewidth, height=5.4cm,
      xlabel={Cores $N$}, ylabel={measured AMAT (cycles)},
      xmode=log, log basis x=2, xtick={2,4,8,16}, xticklabels={2,4,8,16},
      ymin=0, ymax=85, grid=major, grid style={coreGrey!15},
      label style={font=\footnotesize}, tick label style={font=\scriptsize},
      legend style={at={(0.03,0.97)}, anchor=north west, font=\tiny, draw=none, fill=white}]
    \addplot[very thick, mesiBlue, mark=*] table[x=N, y=Am] {data/e2_heavy.dat}; \addlegendentry{MESI, heavy (think 40)}
    \addplot[very thick, moesiOrange, mark=square*] table[x=N, y=Ao] {data/e2_heavy.dat}; \addlegendentry{MOESI, heavy}
    \addplot[thick, dashed, mesiBlue, mark=o] table[x=N, y=Am] {data/e2_light.dat}; \addlegendentry{MESI, light (think 300)}
    \addplot[thick, dashed, moesiOrange, mark=square] table[x=N, y=Ao] {data/e2_light.dat}; \addlegendentry{MOESI, light}
  \end{axis}
\end{tikzpicture}
\caption[AMAT against cores]{Experiment E2: AMAT grows with cores and the MESI--MOESI
gap widens as the bus gets busier. $d=0.5$.}
\label{fig:res-amat}
\end{minipage}\hfill
\begin{minipage}[c]{0.45\linewidth}\centering
\begin{tikzpicture}
  \begin{axis}[ybar, bar width=7pt, width=\linewidth, height=5.4cm,
      xlabel={memory-port occupancy (cycles)},
      ylabel={memory-port utilisation (\%)},
      xtick={0,1,2,3}, xticklabels={5,10,20,40},
      ymin=0, ymax=125, enlarge x limits=0.15, grid=major, grid style={coreGrey!15},
      label style={font=\footnotesize}, tick label style={font=\scriptsize},
      legend style={at={(0.03,0.97)}, anchor=north west, font=\tiny, draw=none, fill=white}]
    \addplot[fill=mesiBlue!70, draw=mesiBlue] table[x=x, y=MUm] {data/e5_memport.dat}; \addlegendentry{MESI}
    \addplot[fill=moesiOrange!75, draw=moesiOrange] table[x=x, y=MUo] {data/e5_memport.dat}; \addlegendentry{MOESI}
    \node[font=\tiny, text=memGreen!70!black] at (axis cs:0,28) {$-12.6\%$};
    \node[font=\tiny, text=memGreen!70!black] at (axis cs:1,45) {$-12.6\%$};
    \node[font=\tiny, text=memGreen!70!black] at (axis cs:2,78) {$-13.6\%$};
    \node[font=\tiny\bfseries, text=memGreen!70!black] at (axis cs:3,109) {$-35.9\%$};
  \end{axis}
\end{tikzpicture}
\caption[Memory-port utilisation]{Experiment E5 (8 cores, think 40, $d=0.5$):
memory-port use, with MOESI's time reduction above each pair.}
\label{fig:res-memory}
\end{minipage}
\end{figure}

\begin{table}[htbp]
\centering
\caption{Experiment E2: AMAT and dirty-read latency (cycles), with the
execution-time reduction for comparison.}
\label{tab:res-amat}
\small
\begin{tabular}{@{}l c c c c c c c c@{}}
\toprule
Load & $N$ & Bus util. & \multicolumn{3}{c}{AMAT} &
  \multicolumn{2}{c}{Dirty-read latency} & Time \\
\cmidrule(lr){4-6}\cmidrule(lr){7-8}
 & & MESI (\%) & MESI & MOESI & cut (\%) & MESI & MOESI & cut (\%) \\
\midrule
light & 2 & 2.8 & 14.4 & 12.8 & 11.2 & 61.1 & 52.0 & 0.5 \\
light & 4 & 7.1 & 22.3 & 18.1 & 18.7 & 61.5 & 52.2 & 1.3 \\
light & 8 & 15.5 & 29.0 & 22.6 & 21.9 & 62.4 & 52.6 & 1.8 \\
light & 16 & 32.5 & 37.6 & 28.7 & 23.7 & 64.7 & 53.8 & 2.4 \\
\midrule
heavy & 2 & 15.7 & 14.5 & 12.9 & 11.2 & 61.8 & 52.3 & 3.0 \\
heavy & 4 & 35.8 & 23.8 & 18.8 & 21.1 & 64.0 & 53.2 & 8.0 \\
heavy & 8 & 69.3 & 35.0 & 25.2 & 27.9 & 71.2 & 56.5 & 12.6 \\
heavy & 16 & 97.0 & 78.6 & 50.2 & 36.0 & 125.5 & 86.6 & 23.8 \\
\bottomrule
\end{tabular}
\end{table}

MOESI lowers AMAT by 11 to 36\,\%, far more than it lowers execution time (0.5 to
23.8\,\%). AMAT averages memory operations only, while execution time also includes
the computation between them, which dominates at light load. And part of the AMAT
cut is Effect~B of \cref{fig:two-effects}: under MOESI the readers after the first
in each epoch are served by the owner in 51 cycles instead of by memory in 100. The
example of \cref{tab:amat-compare} counted Effect~A alone and predicted an AMAT cut
below 10\,\%; the measured cut is larger because of Effect~B. This is why
Chapter~16 kept the two effects apart: only Effect~A is the forced writeback.

Experiment E5 varies how long the memory port is busy per 64-byte access, from 5
cycles (fast memory) to 40 (slow memory) (\cref{fig:res-memory}). With fast memory
the bus is the bottleneck and MOESI saves the same 12.6\,\% as in E1. The memory
port, though, is 17 to 98\,\% busy under MESI across the four settings and 2 to
16\,\% under MOESI, because every forced writeback is also a memory write. At 40
cycles per access MESI saturates the port (98\,\% busy), requests queue at memory,
the bus sits idle half the time (its utilisation falls from 69 to 50\,\%), and
MOESI, which hardly uses the port, finishes 35.9\,\% sooner. This is the
memory-bandwidth argument of Chapter~9 in numbers: the forced writeback costs time
wherever bandwidth is scarce, on the bus or at memory.

% ===============================================================
%  Chapter 20 -- Testing the Dirty Sharing Ratio
% ===============================================================
\chapter{Testing the Dirty Sharing Ratio}
\label{ch:res-dsr}
\index{Dirty Sharing Ratio}%

\begin{figure}[htbp]
\centering
\adjustbox{max width=\linewidth}{%
\begin{forest}
  chapter map,
  [Chapter 20
    [Why not the input knob [$d$ is far from linear]]
    [H3: DSR as predictor [One straight line] [Three cross-checks]]
    [Survival {\&} small caches [$\sigma$ stays near 1] [DSR carries the effect]]
    [H4: when MOESI loses [Never in bytes] [Slow cache-to-cache]]
    [What it means [Published studies] [Design rule] [Verdicts]]
  ]
\end{forest}}
\caption[]{Structure of Chapter~20: why a measured ratio is
needed, how well DSR predicts MOESI's saving, what small caches do, and when MOESI
is slower than MESI.}
\label{fig:ch20-map}
\end{figure}

% ---------------------------------------------------------------
\section{Why the Input Knob Is Not a Predictor}
\label{sec:res-knob}
\sectionmap{[Input knob [d describes the program] [Saving saturates in d]
  [{Same d, different savings}] [Need a measured ratio]]}

The generator's knob $d$, the share of operations that touch dirty-shared data, is
the obvious measure of how much a program shares, so we checked it first
(\cref{fig:res-knob}). The curve is far from straight: the saving rises fast at
small $d$ and levels off near 40\,\%, so the correlation between $d$ and the saving
is only 0.64 (0.93 with $d$ on a log scale). At a given $d$ the saving also depends
strongly on the core count, and at $d=0.1$ single runs range from 25.2 to 42.0\,\%.
Both problems have the cause found in \cref{sec:dsr-patterns}: what matters is how
many transactions the shared accesses produce, not how many accesses there are
(\cref{eq:dsr-mix}).

\begin{figure}[htbp]
\centering
\begin{tikzpicture}
  \begin{axis}[width=10cm, height=5.2cm,
      xlabel={dirty-sharing share of operations, $d$ (log scale)}, ylabel={MOESI saving $S_B$ (\%)},
      xmode=log, xmin=0.0008, xmax=1.2, ymin=0, ymax=48, grid=major, grid style={coreGrey!15},
      label style={font=\small}, tick label style={font=\scriptsize},
      legend style={at={(0.03,0.97)}, anchor=north west, font=\scriptsize, draw=none, fill=white}]
    \addplot[very thick, skyC, mark=*] table[x=d, y=N2] {data/S_vs_d.dat}; \addlegendentry{$N=2$}
    \addplot[very thick, leafC, mark=square*] table[x=d, y=N4] {data/S_vs_d.dat}; \addlegendentry{$N=4$}
    \addplot[very thick, amberC, mark=triangle*] table[x=d, y=N8] {data/S_vs_d.dat}; \addlegendentry{$N=8$}
    \addplot[very thick, roseC, mark=diamond*] table[x=d, y=N16] {data/S_vs_d.dat}; \addlegendentry{$N=16$}
  \end{axis}
\end{tikzpicture}
\caption[Saving against the input knob]{Experiment T1. MOESI's saving rises steeply
at small $d$ and then levels off, and the same $d$ gives different savings at
different core counts.}
\label{fig:res-knob}
\end{figure}

% ---------------------------------------------------------------
\section{H3: DSR Predicts the Traffic Saving}
\label{sec:res-h3}
\sectionmap{[H3 [Scatter of 360 runs] [{One line, $r = 0.990$}] [{Slope falls with N}]
  [Three cross-checks] [Verdict: held]]}
\index{MESI}\index{MOESI}\index{forced writeback}\index{dirty sharing}\index{Dirty Sharing Ratio}\index{survival fraction}\index{verification}%

H3 predicted that the saving obeys $S_B=\sigma\cdot72\,\dsr/\bar b$
(\cref{eq:identity}) exactly, and that with $\sigma$ near 1 and $\bar b$ nearly
constant, DSR measured under MESI predicts it through a straight line
(\cref{eq:predictor}). \Cref{fig:res-dsr} plots the saving of every run with any
dirty sharing (360 runs) against its DSR. Unlike \cref{fig:res-knob}, the points
lie close to one line for all four core counts. \Cref{tab:res-dsr} gives the fits.

\begin{figure}[htbp]
\centering
\begin{tikzpicture}
  \begin{axis}[width=11.5cm, height=6.6cm,
      xlabel={Dirty Sharing Ratio measured under MESI}, ylabel={MOESI traffic saving $S_B$ (\%)},
      xmin=0, xmax=0.52, ymin=-2, ymax=48, grid=major, grid style={coreGrey!15},
      xtick={0,0.1,0.2,0.3,0.4,0.5},
      xticklabel style={/pgf/number format/fixed, /pgf/number format/precision=1},
      label style={font=\small}, tick label style={font=\scriptsize},
      legend style={at={(0.02,0.98)}, anchor=north west, font=\scriptsize, draw=none, fill=white, fill opacity=0.85, text opacity=1}]
    \addplot[only marks, mark=*, mark size=1.4pt, skyC, fill opacity=0.6, draw opacity=0.6] table[x=DSR, y=S] {data/dsr_scatter_N2.dat}; \addlegendentry{$N=2$}
    \addplot[only marks, mark=square*, mark size=1.3pt, leafC, fill opacity=0.6, draw opacity=0.6] table[x=DSR, y=S] {data/dsr_scatter_N4.dat}; \addlegendentry{$N=4$}
    \addplot[only marks, mark=triangle*, mark size=1.6pt, amberC, fill opacity=0.6, draw opacity=0.6] table[x=DSR, y=S] {data/dsr_scatter_N8.dat}; \addlegendentry{$N=8$}
    \addplot[only marks, mark=diamond*, mark size=1.6pt, roseC, fill opacity=0.6, draw opacity=0.6] table[x=DSR, y=S] {data/dsr_scatter_N16.dat}; \addlegendentry{$N=16$}
    \addplot[very thick, black, dashed, domain=0.04:0.52] {92.55*x - 3.63}; \addlegendentry{fit: $S_B = 92.6\,\dsr - 3.6$}
    \addplot[thick, plumC, domain=0:0.52] {90*x}; \addlegendentry{identity, $\bar b = 80$\,B, $\sigma = 1$}
  \end{axis}
\end{tikzpicture}
\caption[H3: saving against DSR]{Experiment T1, all 360 runs with any dirty
sharing. One straight line in DSR fits every core count (Pearson $r=0.990$). The
purple line is \cref{eq:identity} with $\bar b=80$\,B and $\sigma=1$, the values of
the pure dirty-sharing patterns.}
\label{fig:res-dsr}
\end{figure}

\begin{table}[htbp]
\centering
\caption{Straight-line fits $S_B = a\cdot\dsr + b$ (saving in \%), with errors in
points. ``Held out'' fits the line on the other three core counts and predicts this
one.}
\label{tab:res-dsr}
\small
\begin{tabular}{@{}l c c c c c c c@{}}
\toprule
 & Slope $a$ & Intercept $b$ & $r$ & Mean error & Max error & \multicolumn{2}{c}{Held out: mean / max} \\
\midrule
$N=2$  & 97.6 & $-5.9$ & 0.989 & 1.89 & 4.24 & 2.08 & 5.24 \\
$N=4$  & 94.5 & $-4.3$ & 0.990 & 1.85 & 3.37 & 1.86 & 3.62 \\
$N=8$  & 91.7 & $-3.3$ & 0.991 & 1.76 & 3.20 & 1.75 & 3.26 \\
$N=16$ & 89.1 & $-2.5$ & 0.991 & 1.61 & 2.90 & 1.78 & 3.96 \\
\midrule
All    & 92.6 & $-3.6$ & 0.990 & 1.84 & 5.10 & --- & --- \\
\bottomrule
\end{tabular}
\end{table}

One line fitted to all 360 runs predicts the saving with a mean error of 1.84
points; 95\,\% of runs are within 3.56 points and the worst is 5.10 points off. A
separate line per core count does slightly better (mean error 1.61 to 1.89 points,
worst 4.24). The identity makes the slope proportional to $1/\bar b$, and $\bar b$
rises with the core count (\cref{tab:res-h1}), so the slope falls from 97.6 at 2
cores to 89.1 at 16. The fitted slopes are also steeper, with negative intercepts,
than a line of slope $7200/\bar b$ through the origin, because $\bar b$ is not
constant within a core count. Workloads with little dirty sharing have a larger
$\bar b$ (up to 137\,B, since private misses and evictions add bytes without dirty
reads), while workloads dominated by dirty sharing approach the 80\,B of the pure
patterns. The saving therefore rises a little faster than linearly in DSR, and the
best straight line starts below zero.

A predictor has to work on data it was not fitted to, and we tested this three
ways. A line fitted on three core counts predicts the fourth with a mean error of
1.75 to 2.08 points and a worst error of 5.24 (\cref{tab:res-dsr}, last columns). A
per-core-count line fitted on two seeds predicts the third with a mean error of 1.73
to 1.84 points. And the line fitted on the 64\,KB runs of T1 predicts the 72 runs
of T2, with caches from 4 to 1024 blocks, with a mean error of 1.40 points and a
worst of 3.01. In every case the error on unseen data is close to the error on the
fitted data, so the line is not just memorising the runs it was fitted to.

The identity itself held to rounding error in all 468 traffic runs
(\cref{tab:verify-identities}). With $\bar b$ measured alongside DSR and $\sigma$
taken as 1, it predicts the saving with a mean error of 0.31 points and a worst of
1.85, always a slight overestimate. That is much better than the line, but it tells
us little: $72\,\dsr/\bar b$ is exactly the forced-writeback fraction $F$, so this
prediction only says that MOESI removes nearly all forced writebacks ($\sigma$
between 0.83 and 1), which in a byte-counting model is close to true by
construction. The useful result is the straight line, which needs only the two
counters of \cref{fig:dsr-counter}. Whether it carries over to real programs, with
their timing and working sets, needs real traces (Chapter~22).

\begin{verdictbox}{Verdict on H3: held (within this model)}
The identity held exactly in every run. DSR alone, through one straight line,
predicted MOESI's traffic saving to within 1.84 points on average, with the same
accuracy on core counts, traces and cache sizes it was not fitted to. The slope
depends on $\bar b$, so the line should be fitted for each system.
\end{verdictbox}

% ---------------------------------------------------------------
\section{The Survival Fraction and Small Caches}
\label{sec:res-sigma}
\sectionmap{[Survival [sigma between 0.83 and 1] [Small caches cut the saving]
  [{Mainly through DSR, not sigma}] [Deferred write backs appear]]}
\index{MESI}\index{MOESI}\index{Dirty Sharing Ratio}\index{survival fraction}\index{dirty epoch}\index{owner}%

Part~IV expected small caches to reduce MOESI's gain by lowering $\sigma$, since an
owner evicted before the next write pays the writeback late instead of never
(\cref{fig:sigma-tree}). Experiment T2 tests this at 8 cores with caches of 4 to
1024 blocks (\cref{fig:res-sigma,tab:res-sigma}).

\begin{figure}[htbp]
\centering
\begin{bigfig}[1.0]
\begin{tikzpicture}
  \begin{axis}[name=a, width=7.4cm, height=5.6cm,
      title={\sffamily\small (a) Dirty Sharing Ratio},
      xlabel={cache capacity (blocks)}, ylabel={$\dsr$ under MESI},
      xmode=log, log basis x=2, xtick={4,16,64,256,1024}, xticklabels={4,16,64,256,1024},
      ymin=0, ymax=0.5, grid=major, grid style={coreGrey!15},
      yticklabel style={/pgf/number format/fixed, /pgf/number format/precision=1},
      label style={font=\small}, tick label style={font=\scriptsize},
      legend style={at={(1.08,-0.28)}, anchor=north, legend columns=3, font=\scriptsize, draw=none}]
    \addplot[very thick, mesiBlue, mark=*] table[x=cap, y=dsrd0p2] {data/t2_cache.dat}; \addlegendentry{$d=0.2$}
    \addplot[very thick, memGreen, mark=square*] table[x=cap, y=dsrd0p5] {data/t2_cache.dat}; \addlegendentry{$d=0.5$}
    \addplot[very thick, moesiOrange, mark=triangle*] table[x=cap, y=dsrd0p8] {data/t2_cache.dat}; \addlegendentry{$d=0.8$}
  \end{axis}
  \begin{axis}[at={(a.east)}, anchor=west, xshift=1.5cm, width=7.4cm, height=5.6cm,
      title={\sffamily\small (b) Survival fraction $\sigma$},
      xlabel={cache capacity (blocks)}, ylabel={$\sigma$},
      xmode=log, log basis x=2, xtick={4,16,64,256,1024}, xticklabels={4,16,64,256,1024},
      ymin=0, ymax=1.05, grid=major, grid style={coreGrey!15},
      yticklabel style={/pgf/number format/fixed, /pgf/number format/precision=1},
      label style={font=\small}, tick label style={font=\scriptsize}]
    \addplot[very thick, mesiBlue, mark=*] table[x=cap, y=sigd0p2] {data/t2_cache.dat};
    \addplot[very thick, memGreen, mark=square*] table[x=cap, y=sigd0p5] {data/t2_cache.dat};
    \addplot[very thick, moesiOrange, mark=triangle*] table[x=cap, y=sigd0p8] {data/t2_cache.dat};
  \end{axis}
\end{tikzpicture}
\end{bigfig}
\caption[DSR and survival fraction against cache size]{Experiment T2, 8 cores.
(a) Small caches lower DSR sharply: owners are evicted before other cores read their
data. (b) The survival fraction falls much less, from about 1 to about 0.89.}
\label{fig:res-sigma}
\end{figure}

\begin{table}[htbp]
\centering
\caption{Experiment T2 at $d = 0.5$: what small caches change. Eviction writebacks
are totals per run.}
\label{tab:res-sigma}
\small
\begin{tabular}{@{}r c c c c c c@{}}
\toprule
Capacity & MESI hit rate & $\dsr$ & $\sigma$ & MOESI saving & \multicolumn{2}{c}{Eviction writebacks} \\
\cmidrule(lr){6-7}
(blocks) & (\%) & & & (\%) & MESI & MOESI \\
\midrule
4 & 22.1 & 0.124 & 0.895 & 8.2 & 12,292 & 12,976 \\
8 & 24.5 & 0.190 & 0.890 & 13.2 & 8,146 & 9,286 \\
16 & 29.4 & 0.248 & 0.893 & 18.1 & 4,630 & 6,068 \\
32 & 38.6 & 0.286 & 0.939 & 22.4 & 3,472 & 4,337 \\
64 & 52.3 & 0.336 & 0.994 & 28.1 & 3,027 & 3,094 \\
128 & 67.0 & 0.422 & 0.999 & 36.8 & 468 & 468 \\
256 & 68.0 & 0.429 & 0.999 & 37.8 & 0 & 0 \\
1024 & 68.0 & 0.429 & 0.999 & 37.8 & 0 & 0 \\
\bottomrule
\end{tabular}
\end{table}

As the cache shrinks from 256 to 4 blocks, MOESI's saving falls from 37.8 to
8.2\,\%. Some of this is $\sigma$, which falls from 0.999 to 0.895: about one
dirty-read epoch in ten now ends with the owner evicted, and MOESI makes more
eviction writebacks than MESI (12,976 against 12,292 at 4 blocks). Those
writebacks were deferred, not removed, as \cref{sec:own-moesi} warned. Most of the
fall, though, comes from DSR, which drops from 0.429 to 0.124. With tiny caches the
writer's copy is often evicted, and written back by both protocols, before another
core reads it, so the reader finds the block in memory and neither protocol has a
dirty read. Capacity misses also add many ordinary transactions that dilute the
dirty reads. Because the effect of cache size shows up in DSR itself, the DSR line
still predicts the saving when caches shrink, which is why the third cross-check of
\cref{sec:res-h3} worked. In T1, with 64\,KB caches, $\sigma$ ranged from 0.83 to 1
with a mean of 0.97. Its lowest values came at the smallest $d$, where a run has few
dirty epochs and the writebacks charged at the end weigh more.

% ---------------------------------------------------------------
\section{H4: When Does MOESI Lose?}
\label{sec:res-h4}
\sectionmap{[H4 [Never in bytes] [Never in time with report latencies]
  [{Loses when cache-to-cache is slow}] [Effect B] [Verdict: held]]}
\index{MESI}\index{MOESI}\index{Dirty Sharing Ratio}\index{survival fraction}\index{dirty epoch}\index{cache-to-cache transfer}\index{owner}\index{Effect A and Effect B}%

H4 predicted that MOESI never moves more bytes than MESI and can take more time only
through Effect~B (a cache-to-cache read slower than a memory read) or owner
contention, and that small caches lower the gain without reversing it. MOESI moved
no more bytes than MESI in any of the 468 traffic runs or 288 timing runs,
including the smallest caches. With $T_{\mathrm{coh}}=60$ it was also faster than
MESI in all 228 timing runs of E1, E2, E3 and E5. In experiment E3
(\cref{fig:res-h4}a) small caches shrink the time gain as they shrank the traffic
gain, but at the worst point (4 blocks, $d=0.2$) MOESI is still 0.3\,\% faster.

\begin{figure}[htbp]
\centering
\begin{minipage}[t]{0.47\linewidth}
\centering
\adjustbox{max width=\linewidth}{%
\begin{tikzpicture}
  \begin{axis}[width=7.6cm, height=6.4cm,
      title={\sffamily\small (a) Cache capacity (E3)},
      xlabel={cache capacity (blocks)}, ylabel={execution-time reduction (\%)},
      xmode=log, log basis x=2, xtick={4,16,64,256,1024}, xticklabels={4,16,64,256,1024},
      xmin=3, xmax=1400, ymin=-1, ymax=20, grid=major, grid style={coreGrey!15},
      label style={font=\small}, tick label style={font=\scriptsize},
      legend style={at={(0.97,0.03)}, anchor=south east, font=\scriptsize, draw=none, fill=white}]
    \addplot[very thick, mesiBlue, mark=*] table[x=cap, y=redd0p2] {data/e3_cache.dat}; \addlegendentry{$d=0.2$}
    \addplot[very thick, memGreen, mark=square*] table[x=cap, y=redd0p5] {data/e3_cache.dat}; \addlegendentry{$d=0.5$}
    \addplot[very thick, moesiOrange, mark=triangle*] table[x=cap, y=redd0p8] {data/e3_cache.dat}; \addlegendentry{$d=0.8$}
    \addplot[thick, dashed, coreGrey, forget plot] coordinates {(3,0) (1400,0)};
  \end{axis}
\end{tikzpicture}}
\end{minipage}\hfill
\begin{minipage}[t]{0.51\linewidth}
\centering
\adjustbox{max width=\linewidth}{%
\begin{tikzpicture}[font=\sffamily\scriptsize,
    cell/.style={draw=white, line width=1pt, minimum width=1.6cm, minimum height=0.75cm, font=\sffamily\small}]
  \node[font=\sffamily\small] at (2.4,1.45) {(b) Latency sensitivity (E4)};
  \node[font=\sffamily\scriptsize] at (2.4,1.05) {owner-port occupancy (cycles)};
  % generated by make_report_data.py -- do not edit by hand
\node[cell, fill=alertRed!26, text=black] at (0.0,0.00) {-4.4};
\node[cell, fill=alertRed!26, text=black] at (1.6,0.00) {-4.4};
\node[cell, fill=alertRed!28, text=black] at (3.2,0.00) {-5.1};
\node[cell, fill=alertRed!30, text=black] at (4.8,0.00) {-5.7};
\node[cell, fill=skyC!12, text=black] at (0.0,-0.75) {+0.1};
\node[cell, fill=skyC!12, text=black] at (1.6,-0.75) {+0.1};
\node[cell, fill=alertRed!12, text=black] at (3.2,-0.75) {0.0};
\node[cell, fill=alertRed!14, text=black] at (4.8,-0.75) {-0.5};
\node[cell, fill=skyC!32, text=black] at (0.0,-1.50) {+6.4};
\node[cell, fill=skyC!32, text=black] at (1.6,-1.50) {+6.4};
\node[cell, fill=skyC!32, text=black] at (3.2,-1.50) {+6.3};
\node[cell, fill=skyC!30, text=black] at (4.8,-1.50) {+5.6};
\node[cell, fill=skyC!52, text=white] at (0.0,-2.25) {+12.6};
\node[cell, fill=skyC!52, text=white] at (1.6,-2.25) {+12.6};
\node[cell, fill=skyC!53, text=white] at (3.2,-2.25) {+13.0};
\node[cell, fill=skyC!48, text=white] at (4.8,-2.25) {+11.3};
\node[cell, fill=skyC!70, text=white] at (0.0,-3.00) {+18.3};
\node[cell, fill=skyC!70, text=white] at (1.6,-3.00) {+18.3};
\node[cell, fill=skyC!70, text=white] at (3.2,-3.00) {+18.1};
\node[cell, fill=skyC!65, text=white] at (4.8,-3.00) {+16.7};
\node[font=\sffamily\scriptsize] at (0.0,0.62) {5};
\node[font=\sffamily\scriptsize] at (1.6,0.62) {10};
\node[font=\sffamily\scriptsize] at (3.2,0.62) {20};
\node[font=\sffamily\scriptsize] at (4.8,0.62) {40};
\node[font=\sffamily\scriptsize, anchor=east] at (-0.9,0.00) {200 (c2c 191)};
\node[font=\sffamily\scriptsize, anchor=east] at (-0.9,-0.75) {150 (c2c 141)};
\node[font=\sffamily\scriptsize, anchor=east] at (-0.9,-1.50) {100 (c2c 91)};
\node[font=\sffamily\scriptsize, anchor=east] at (-0.9,-2.25) {60 (c2c 51)};
\node[font=\sffamily\scriptsize, anchor=east] at (-0.9,-3.00) {30 (c2c 21)};

  \node[font=\sffamily\scriptsize, rotate=90] at (-3.55,-1.5) {$T_{\mathrm{coh}}$ (c2c latency)};
  \draw[very thick, memGreen] (-0.82,-2.63) rectangle (5.62,-1.87);
  \node[font=\sffamily\tiny, text=coreGrey, align=center] at (2.4,-4.05) {cells: MOESI execution-time reduction (\%);\\blue = MOESI faster, red = MOESI slower;\\green frame = the report's $T_{\mathrm{coh}} = 60$};
\end{tikzpicture}}
\end{minipage}
\caption[H4: when MOESI loses]{(a) Experiment E3: small caches reduce MOESI's time
gain but never reverse it. (b) Experiment E4: MOESI becomes slower only when the
cache-to-cache latency is well above the 100-cycle memory latency. Rows:
$T_{\mathrm{coh}}$ and the resulting cache-to-cache latency; columns: owner-port
occupancy. 8 cores, $d=0.5$, means of three seeds.}
\label{fig:res-h4}
\end{figure}

Experiment E4 raises the cost of a cache-to-cache transfer (\cref{fig:res-h4}b).
MOESI's advantage shrinks steadily as $T_{\mathrm{coh}}$ rises, reaches zero near
$T_{\mathrm{coh}}=150$ (a 141-cycle cache-to-cache read against 100 cycles from
memory), and becomes a loss of up to 5.7\,\% on average and 5.9\,\% in the worst run
at $T_{\mathrm{coh}}=200$. MOESI was slower in 16 of the 60 E4 runs, all at
$T_{\mathrm{coh}}\ge150$. The time the owner's cache spends supplying a block
matters much less: raising it from 5 to 40 cycles lowers each row by at most 1.6
points. \Cref{tab:res-effectb} shows the mechanism is Effect~B. Under MOESI the
owner answers every read in an epoch, so it supplies about 45\,\% more reads than
MESI's dirty reads, and MESI served those extra reads from memory. With a 191-cycle
cache-to-cache read, each of them is 91 cycles slower than a memory read would have
been. At $T_{\mathrm{coh}}=200$ MOESI's dirty reads are still faster than MESI's
(194.8 against 207.0 cycles), yet the run as a whole is slower, so a study that
measured only the read MOESI was designed to improve would miss the loss.

\begin{table}[htbp]
\centering
\caption{Experiment E4, means over the four owner-port occupancies. ``Reads from a
cache'' counts read misses supplied by another cache (under MESI, the dirty reads).}
\label{tab:res-effectb}
\small
\begin{tabular}{@{}c c c c c c c c@{}}
\toprule
$T_{\mathrm{coh}}$ & c2c & \multicolumn{3}{c}{Reads from a cache, per run} &
  \multicolumn{2}{c}{Their latency} & Time saved \\
\cmidrule(lr){3-5}\cmidrule(lr){6-7}
 & & MESI & MOESI & ratio & MESI & MOESI & (\%) \\
\midrule
30 & 21 & 6,123 & 8,628 & 1.41 & 43.9 & 29.2 & $+17.8$ \\
60 & 51 & 5,990 & 8,625 & 1.44 & 71.5 & 57.6 & $+12.4$ \\
100 & 91 & 5,882 & 8,611 & 1.46 & 109.7 & 96.4 & $+6.2$ \\
150 & 141 & 5,759 & 8,596 & 1.49 & 158.1 & 145.4 & $-0.1$ \\
200 & 191 & 5,677 & 8,594 & 1.51 & 207.0 & 194.8 & $-4.9$ \\
\bottomrule
\end{tabular}
\end{table}

\begin{verdictbox}[amberC]{Verdict on H4: held, with one correction}
MOESI never moved more bytes than MESI and was slower only when a cache-to-cache
transfer was much slower than a memory read, which is Effect~B as predicted. Owner
contention mattered little. Small caches did reduce the gain without reversing it,
but mainly by lowering DSR, not $\sigma$ (\cref{sec:res-sigma}).
\end{verdictbox}

% ---------------------------------------------------------------
\section{What the Results Mean}
\label{sec:res-meaning}
\label{sec:res-glance}
\sectionmap{[Meaning [Agrees with the published evaluations] [Explains why they differ]
  [A design rule in four questions] [Four verdicts]]}
\index{MESI}\index{MOESI}\index{Owned state}\index{bus utilisation}%

Chapter~14 found that published evaluations disagree on how much MOESI helps and
that none isolates \stO\ (\cref{tab:evaluations}). Our controlled results do not
replace those studies, which use real programs and detailed simulators, but they
give one consistent explanation of what the studies report (\cref{tab:res-published}).
The third row is the most telling. A protocol can remove most memory writes (MOESI
removed 83 to 99.8\,\% of them for $d\ge0.1$, \cref{tab:res-traffic}) and still
change run time by about one per cent, because a system with spare bandwidth does
not wait for those writes. Measured on such a system MOESI and MESI look equal; on a
saturated one MOESI cuts run time by up to 28\,\%. Both measurements are right, and
H2 says which one a given study will see.

\begin{table}[htbp]
\centering
\caption{Published findings (as summarised in Chapters 11 and 14) and the result
here that accounts for each. The right-hand column is an explanation consistent
with our model, not a reproduction of those studies.}
\label{tab:res-published}
\small
\begin{tabular}{@{}L{3.5cm} L{4.2cm} L{5.0cm}@{}}
\toprule
Study & Published finding & Corresponding result here \\
\midrule
Archibald {\&} Baer~\cite[p.~295]{archibald1986evaluation} & the owner-based
  Berkeley protocol surpasses Illinois at high sharing, because it does not
  update memory on a miss to a modified block & the saving rises with DSR
  (\cref{fig:res-dsr}) and is exactly the removed writebacks
  (\cref{eq:identity}) \\
Fagundes {\&} Souza~\cite{fagundes2024evaluation} & MOESI's advantage over
  MESI shrinks as the last-level cache grows & MOESI gains less when memory-side
  reads are cheap or uncongested relative to cache-to-cache reads (E4,
  \cref{fig:res-h4}; E5, \cref{fig:res-memory}) \\
Loughlin et al.~\cite{loughlin2022moesiprime} & removing downgrade writebacks
  cuts DRAM row activations but changes run time by only $-0.51$ to
  $+0.61\,\%$ & saved bytes become time only under load: 0.3 to 1.8\,\% below
  16\,\% bus utilisation (\cref{tab:res-h2}) \\
\bottomrule
\end{tabular}
\end{table}

The four verdicts combine into a rule a designer could apply before building MOESI,
using only quantities measurable on a MESI system (\cref{fig:res-rule}). The first
question needs the two counters of \cref{fig:dsr-counter}, and its answer also gives
the expected traffic saving through the line of \cref{tab:res-dsr}. The second needs
the utilisation of the bus and memory port on the existing system, and the third
needs two latencies the designer already knows. The leaves give numbers only where
Part~VI measured a range, so they share its limitations (Chapter~21).
\Cref{tab:res-glance} collects the four verdicts; Chapter~22 turns them into
answers to the research questions.

\begin{figure}[htbp]
\centering
\adjustbox{max width=0.95\linewidth}{%
\begin{forest}
  decision tree,
  [{Workload running\\on a MESI system}
    [{Is DSR above zero?\\(two counters)}
      [{Is the bus or memory\\port more than a third busy?}, el=yes
        [{Is a cache-to-cache read\\faster than a memory read?}, el=yes
          [{MOESI saves traffic {\it and} time:\\time 5--28\,\% here,\\growing with utilisation}, el=yes]
          [{Risk: MOESI may be\\slower (Effect B),\\up to 6\,\% here}, el=no]]
        [{MOESI saves traffic\\and memory writes,\\but little time ($\le 4\,\%$ here)}, el=no]]
      [{No forced write backs:\\MOESI cannot help}, el=no]]
  ]
\end{forest}}
\caption[A design rule]{A design rule from the four verdicts. Every question can be
answered on the existing MESI system; ranges are those measured in this Part.}
\label{fig:res-rule}
\end{figure}

\begin{table}[htbp]
\centering
\caption{The four hypotheses of Part~IV against the measurements of Part~VI.}
\label{tab:res-glance}
\small
\begin{tabular}{@{}L{1.3cm} L{3.7cm} L{5.3cm} L{2.2cm}@{}}
\toprule
 & Predicted & Measured & Verdict \\
\midrule
H1 (RQ1) & $F$ falls as cores are added & $F$ falls at every mix; 43.6\,\% at 2 cores
  to 33.4\,\% at 16 for $d \ge 0.1$ & held \\
H2 (RQ2) & time gain follows utilisation; limit = bus-busy ratio &
  1\,\% of the byte saving becomes time at 2.5\,\% utilisation, 95\,\% at
  saturation; limit reached & held \\
H3 (RQ3) & identity exact; DSR predicts the saving & identity exact in all runs;
  one line, mean error 1.84 points, holds on unseen data & held (in this model) \\
H4 (RQ4) & never more bytes; slower only via Effect~B & never more bytes; slower
  only if c2c $\gg$ memory; small caches act through DSR & held, corrected \\
\bottomrule
\end{tabular}
\end{table}
